% =====================================================================
%  Condensed portfolio: one project per page.
%
%  HOW TO EDIT
%  - On Overleaf: upload this whole "portfolio-pdf" folder (portfolio.tex
%    plus the images folder), then press Recompile.
%  - Each project is one \project{...} block at the bottom of this file.
%    Change the words between the braces. Keep the braces themselves.
%  - Every \item is ONE short sentence. Three items per column at most.
%  - To add a project, copy a whole \project block and change it.
%    To remove one, delete its block.
%    \projectinprogress is the same thing without pictures.
%  - Write \% for a percent sign and \& for an ampersand.
%
%  When you push a change to GitHub, the PDF on the website is rebuilt
%  for you (see .github/workflows/portfolio-pdf.yml).
% =====================================================================
\documentclass[11pt]{article}
\usepackage[landscape, letterpaper, margin=0.6in, top=0.55in, bottom=0.55in]{geometry}
\usepackage[T1]{fontenc}
\usepackage[utf8]{inputenc}
\usepackage{helvet}
\renewcommand{\familydefault}{\sfdefault}
\usepackage{graphicx}
\usepackage{xcolor}
\usepackage{enumitem}
\usepackage[hidelinks]{hyperref}
\graphicspath{{images/}}
\pagestyle{empty}
\setlength{\parindent}{0pt}

% ---- Your details: change them here once ----
\newcommand{\myname}{Ariadne Colina-Valeri}
\newcommand{\myemail}{ariadne.colina\_valeri@tufts.edu}
\newcommand{\myemailurl}{mailto:ariadne.colina_valeri@tufts.edu}
\newcommand{\mysite}{acolin03.github.io/Portfolio\_Ariadne.Colina\_Valeri}
\newcommand{\mysiteurl}{https://acolin03.github.io/Portfolio_Ariadne.Colina_Valeri/}

% ---- Colors (the same blue, lilac and sea green as the website) ----
\definecolor{ink}{HTML}{1F2A44}
\definecolor{sub}{HTML}{5B6781}
\definecolor{blue}{HTML}{4C63A0}
\definecolor{lilac}{HTML}{CDBBE6}
\definecolor{sea}{HTML}{A8D6C3}
\color{ink}

\setlist[itemize]{leftmargin=1.1em, itemsep=3pt, topsep=4pt, parsep=0pt, label={\color{lilac}\rule[0.15ex]{0.45em}{0.45em}}}

% A column heading: PROBLEM, SOLUTION, RESULT
\newcommand{\colhead}[1]{\raggedright{\small\bfseries\color{blue}\MakeUppercase{#1}}\par}

% One page per project.
%   #1 title            #2 course and date      #3 skills, separated by \textbar
%   #4 first picture    #5 its caption          #6 second picture    #7 its caption
%   #8 problem items    #9 solution and result (use \solution and \result inside)
\newcommand{\project}[9]{%
  \clearpage
  {\small\color{sub}\myname \hfill \href{\myemailurl}{\myemail} \quad\textbar\quad \href{\mysiteurl}{\mysite}}\par
  \vspace{0.16in}
  {\footnotesize\bfseries\color{blue}\MakeUppercase{#2}}\par
  \vspace{0.09in}
  {\fontsize{26}{30}\selectfont\bfseries #1}\par
  \vspace{0.06in}
  {\color{lilac}\rule{0.55in}{3.5pt}}{\color{sea}\rule{0.55in}{3.5pt}}\par
  \vspace{0.07in}
  {\small\color{sub}#3}\par
  \vspace{0.2in}
  \begin{minipage}[t]{0.485\textwidth}\centering
    \raisebox{-\height}{\includegraphics[width=\linewidth, height=3.1in, keepaspectratio]{#4}}\par
    \vspace{3pt}{\footnotesize\color{sub}#5}
  \end{minipage}\hfill
  \begin{minipage}[t]{0.485\textwidth}\centering
    \raisebox{-\height}{\includegraphics[width=\linewidth, height=3.1in, keepaspectratio]{#6}}\par
    \vspace{3pt}{\footnotesize\color{sub}#7}
  \end{minipage}\par
  \vspace{0.22in}
  \begin{minipage}[t]{0.31\textwidth}
    \colhead{Problem}
    \begin{itemize}#8\end{itemize}
  \end{minipage}\hfill
  #9
}
\newcommand{\solution}[1]{\begin{minipage}[t]{0.31\textwidth}\colhead{Solution}\begin{itemize}#1\end{itemize}\end{minipage}\hfill}
\newcommand{\result}[1]{\begin{minipage}[t]{0.31\textwidth}\colhead{Result}\begin{itemize}#1\end{itemize}\end{minipage}}

% A project that has no pictures yet (something still being built).
%   #1 title   #2 course and date   #3 skills   #4 problem items   #5 \sofar{...} and \nextsteps{...}
\newcommand{\projectinprogress}[5]{%
  \clearpage
  {\small\color{sub}\myname \hfill \href{\myemailurl}{\myemail} \quad\textbar\quad \href{\mysiteurl}{\mysite}}\par
  \vspace{0.16in}
  {\footnotesize\bfseries\color{blue}\MakeUppercase{#2}}\par
  \vspace{0.09in}
  {\fontsize{26}{30}\selectfont\bfseries #1}\par
  \vspace{0.06in}
  {\color{lilac}\rule{0.55in}{3.5pt}}{\color{sea}\rule{0.55in}{3.5pt}}\par
  \vspace{0.07in}
  {\small\color{sub}#3}\par
  \vspace{0.45in}
  \begin{minipage}[t]{0.31\textwidth}
    \colhead{Problem}
    \begin{itemize}#4\end{itemize}
  \end{minipage}\hfill
  #5
}
\newcommand{\sofar}[1]{\begin{minipage}[t]{0.31\textwidth}\colhead{So far}\begin{itemize}#1\end{itemize}\end{minipage}\hfill}
\newcommand{\nextsteps}[1]{\begin{minipage}[t]{0.31\textwidth}\colhead{Next}\begin{itemize}#1\end{itemize}\end{minipage}}

\begin{document}

% =====================================================================
%  PROJECTS
% =====================================================================

\projectinprogress
  {Autonomous ESP32 Robot}
  {In progress \textbar{} EE 31 Junior Design \textbar{} Fall 2026}
  {ESP32 \textbar{} Arduino \textbar{} IR LED \& photodetector \textbar{} PCB design \textbar{} Onshape}
  {
    \item The robot has to follow a colored lane and avoid obstacles without human input.
    \item It also has to receive commands from a partner team's robot.
    \item Every decision it makes depends on sensor readings it can trust.
  }
  {
    \sofar{
      \item Programmed a 7-state state machine on the ESP32, advanced by a pushbutton interrupt or a websocket command.
      \item Selected the IR LED and photodetector resistor values and subtract the room's background light from each reading.
      \item Developing detection logic that filters noise in the color readings before the robot decides which lane color it is on.
    }
    \nextsteps{
      \item Designing the robot's custom ESP32 PCB.
      \item Designing its outer body in Onshape.
    }
  }

\project
  {DJ Board on an STM32}
  {EE 14 Embedded Systems \textbar{} Jan. 2026 to May 2026}
  {STM32L432KC \textbar{} C \textbar{} ADC \& DAC \textbar{} Oscilloscope}
  {dj-circuit.png}{The board: microphone, speaker and effect buttons}
  {dj-scope.png}{Oscilloscope captures of the effect outputs}
  {
    \item A real-time audio effects board records sound and remixes it with 7 effect buttons.
    \item Every effect depends on a clean recording.
    \item The first playback came out with a high-pitched hiss and clicks.
  }
  {
    \solution{
      \item Wrote a timer-driven state machine so every sample is recorded and played at a steady rate.
      \item Gave each microphone reading more time to finish so no sample was cut short.
      \item Held the speaker at a steady resting level between clips and faded each clip in and out.
    }
    \result{
      \item An oscilloscope confirmed that the playback matched the original recording.
      \item All 7 effects run using only 64KB of on-chip memory.
    }
  }

\project
  {FPGA Audio Synthesizer}
  {ES 4 Digital Logic Circuits \textbar{} Sept. 2025 to Dec. 2025}
  {SystemVerilog \textbar{} FPGA \textbar{} PWM}
  {fpga-board.png}{The full board}
  {fpga-octaves.png}{The switch that selects the octave}
  {
    \item Each octave doubles a note's frequency.
    \item The first design calculated every note by dividing on the chip.
    \item Division needed over 250\% of the FPGA's space.
  }
  {
    \solution{
      \item Hand-calculated one octave and used bit shifts for the rest, since dividing by 2 is the same as shifting bits.
      \item Mixed the notes with PWM after testing showed that chords distorted.
    }
    \result{
      \item An 8-octave synthesizer, from 16Hz to 8kHz, for a MIDI keyboard.
      \item 6 notes play clearly at once.
    }
  }

\project
  {CTGAN Design Discovery for RRAM}
  {Tufts Emerging Circuits and Systems Lab \textbar{} Sept. 2025 to April 2026}
  {Python \textbar{} PyTorch \textbar{} pandas \textbar{} CUDA \textbar{} CTGAN \textbar{} PINN}
  {rram-pareto.png}{Final results with 5,000 training samples}
  {rram-evolution.png}{The same search with 500 samples (left) and 2,500 (right)}
  {
    \item A published AI framework designs RRAM memory devices.
    \item Its model averaged the training data, which smoothed away extreme designs.
  }
  {
    \solution{
      \item Replaced the model with a CTGAN, which searches those extremes more aggressively.
      \item Scaled the training data 10 times, to 5,000 samples, so the CTGAN could learn reliably.
      \item Used the framework's physics model as a referee to reject impossible designs.
    }
    \result{
      \item 25\% higher endurance than the published result.
      \item 11\% lower switching energy.
      \item Slightly slower switching: 6.57ns vs 4.07ns.
    }
  }

\project
  {UAV Pitch Control and RF Signal Restoration}
  {EE 23 Linear Systems \textbar{} Jan. 2026 to May 2026}
  {MATLAB \textbar{} PID control \textbar{} Wiener filtering}
  {uav-controllers.png}{P, PI, PD and PID controllers on the same UAV}
  {uav-wiener.png}{Wiener filter error (left) and the restored RF signal (right)}
  {
    \item Keeping a UAV level in pitch is the first step toward collecting clear RF signals from the air.
    \item Even a stable UAV shakes slightly and blurs the signals it receives.
  }
  {
    \solution{
      \item Showed with Routh-Hurwitz analysis that P and PI controllers can never stabilize the UAV.
      \item Searched 15,625 gain combinations for the best PID settings.
      \item Restored a blurred, noisy RF signal with a Wiener filter.
    }
    \result{
      \item 52ms rise time, 2.6\% overshoot and 500ms settling time, within real UAV specs.
      \item 77\% lower error than a direct inverse (MSE 2.272 to 0.521).
    }
  }

\end{document}
